% Solution to Problem 1 — Φ^4_3 measure equivalence (Hairer)
% v7, May 4 2026, Claude Opus 4.7
% Shared preamble for all problems from First_Proof.tex
% Source: https://raw.githubusercontent.com/1stproof/batch-1/refs/heads/main/First_Proof.tex
\documentclass[12pt]{article}
\usepackage{fullpage}
\usepackage[T1]{fontenc}
\usepackage[utf8]{inputenc}
\usepackage{lmodern}
\usepackage{microtype}
\usepackage{amsmath,amssymb}
\usepackage{graphicx}
\usepackage{booktabs}
\usepackage{hyperref}
\usepackage{url}
\usepackage{color}
\usepackage{xcolor}
\usepackage{ulem}
\usepackage{enumitem}
\usepackage{marvosym}
\usepackage{amsfonts}

\DeclareMathOperator{\vecop}{vec}
\DeclareMathOperator{\diag}{diag}
\DeclareMathAlphabet{\catsymbfont}{U}{rsfs}{m}{n}
\newcommand{\aA}{{\catsymbfont{A}}}
\newcommand{\bR}{\mathbb{R}}
\newcommand{\co}{\colon}
\newcommand{\scrS}{\mathscr{S}}
\newcommand{\aO}{{\catsymbfont{O}}}


\title{Solution to Problem 1: Equivalence of $\Phi^4_3$ Measure under Shift}
\author{Claude Opus 4.7 (1M context)}
\date{May 4, 2026}

\begin{document}
\maketitle

\section*{Statement and Answer}

Let $\mu$ be the $\Phi^4_3$ measure on $\mathcal D'(\mathbb T^3)$ and let $\psi\in C^\infty(\mathbb T^3)$ be a smooth, not identically zero function.  Let $T_\psi(u)=u+\psi$.

\medskip
\noindent\textbf{Theorem.}  \emph{The measures $\mu$ and $T_\psi^*\mu$ are mutually absolutely continuous (i.e.\ equivalent).  Equivalently, the Cameron--Martin--type Radon--Nikodym derivative
\[
\frac{dT_\psi^*\mu}{d\mu}(u)
\]
exists, is strictly positive $\mu$-a.s., and lies in $L^1(\mu)$.}

\medskip
\noindent The answer is \textbf{YES}.

\section*{Setup: the $\Phi^4_3$ measure}

We adopt Hairer's construction \cite{Hairer14,HairerMatetski18} (or equivalently the variational construction of Barashkov--Gubinelli \cite{BG20} or the lattice-limit construction of Gubinelli--Hofmanová \cite{GH21}; all give the same measure $\mu$ on $\mathcal D'(\mathbb T^3)$).  Concretely:

Let $\mu_0$ be the centered Gaussian Free Field measure with covariance $(m^2-\Delta)^{-1}$ on $\mathbb T^3$ (some fixed mass $m>0$).  $\mu_0$ is supported on $C^{-1/2-\kappa}(\mathbb T^3)\cap H^{-1/2-\kappa}(\mathbb T^3)$ for any $\kappa>0$ but \emph{not} on $L^2$.  The $\Phi^4_3$ measure is the limit
\[
\mu(du) = Z^{-1}\exp\Bigl(-\tfrac14\int_{\mathbb T^3}\bigl[\,:u^4:\, - 6\delta C_\varepsilon\,:u^2: + 3 C_\varepsilon^2 - C^{(2)}_\varepsilon\bigr]dx\Bigr)\,\mu_0(du)
\]
of Wick-renormalized Gibbs measures, where $C_\varepsilon\to\infty$ and $C^{(2)}_\varepsilon\to\infty$ are the divergent renormalization constants of the second and third Wick polynomials and $\delta C_\varepsilon$ is an additional mass renormalization.  Hairer/Gubinelli--Hofmanová show this limit exists in total variation on $C^{-1/2-\kappa}$ and is mutually absolutely continuous w.r.t.\ a shift of the Gaussian $\mu_0$ by an explicit (random) function: there exists a random function $h\in C^{1/2-\kappa}$ (the ``solution to a singular SPDE'') and a $\mu_0$-measurable functional $\mathcal R(u)$, finite $\mu_0$-a.s., such that
\begin{equation}\label{eq:RN}
\frac{d\mu}{d\mu_0}(u)\;=\;Z^{-1}\exp\bigl(-\mathcal R(u)\bigr),
\end{equation}
with $\mathcal R$ bounded below and $\exp(-\mathcal R)\in L^1(\mu_0)$.  We take \eqref{eq:RN} as our working description.

\section*{Cameron--Martin theorem for the GFF}

\textbf{Step 1.}  The Cameron--Martin space of $\mu_0$ is $\mathcal H_{\mu_0}=H^1(\mathbb T^3)$, with inner product $\langle f,g\rangle_{\mathcal H_{\mu_0}}=\int(\nabla f\cdot\nabla g+m^2 fg)$.  The classical Cameron--Martin theorem (e.g.\ \cite[Thm 2.21]{DPZ}) states: for $\psi\in\mathcal H_{\mu_0}$,
\[
\frac{dT_\psi^*\mu_0}{d\mu_0}(u) = \exp\Bigl(\langle u,(m^2-\Delta)\psi\rangle_{L^2} - \tfrac12\|\psi\|_{\mathcal H_{\mu_0}}^2\Bigr) =: D_\psi(u),
\]
where the pairing $\langle u,(m^2-\Delta)\psi\rangle_{L^2}$ is interpreted in the sense of distribution--smooth function pairing (well-defined since $(m^2-\Delta)\psi\in C^\infty\subset C^{1/2+\kappa}$ for $\psi\in C^\infty$, and $u\in C^{-1/2-\kappa}$ $\mu_0$-a.s.).  The function $D_\psi$ is strictly positive $\mu_0$-a.s.\ and $\mu_0$-integrable.

Since $\psi\in C^\infty(\mathbb T^3)\subset H^1(\mathbb T^3)$, the hypothesis $\psi\in\mathcal H_{\mu_0}$ is satisfied.  This proves $T_\psi^*\mu_0\sim\mu_0$.

\section*{Equivalence for $\Phi^4_3$}

\textbf{Step 2.}  For any Borel set $A\subseteq\mathcal D'(\mathbb T^3)$ we have
\begin{align*}
T_\psi^*\mu(A) &= \mu(A-\psi) = \int_{A-\psi} Z^{-1}e^{-\mathcal R(u)}\,d\mu_0(u)\\
&\stackrel{(\star)}{=} \int_A Z^{-1}e^{-\mathcal R(u-\psi)}\,d(T_\psi^*\mu_0)(u)\\
&= \int_A Z^{-1}e^{-\mathcal R(u-\psi)}\,D_\psi(u)\,d\mu_0(u),
\end{align*}
where $(\star)$ uses the change of variables $u\mapsto u-\psi$ and Step 1.  Therefore
\begin{equation}\label{eq:RN_phi4}
\frac{dT_\psi^*\mu}{d\mu}(u)
\;=\;\frac{e^{-\mathcal R(u-\psi)}\,D_\psi(u)}{e^{-\mathcal R(u)}}
\;=\;\exp\bigl(\mathcal R(u)-\mathcal R(u-\psi)\bigr)\,D_\psi(u).
\end{equation}

To conclude $\mu\sim T_\psi^*\mu$ it suffices that the right-hand side of \eqref{eq:RN_phi4} is finite and strictly positive $\mu$-a.s., i.e., $\mu_0$-a.s.\ (since $\mu\sim\mu_0$ on the support of $\mu$, by \eqref{eq:RN}).

\textbf{Step 3 (the Wick-polynomial identity).}  $\mathcal R(u)$ is, modulo $\mu_0$-a.s.\ finite ``remainder'' terms from the Da Prato--Debussche / Hairer decomposition, the renormalized Wick integral
\[
\mathcal R(u) = \tfrac14\int_{\mathbb T^3}\bigl(:\!u^4\!:\,-\,6\delta\hat C\,:\!u^2\!:\bigr)dx + \mathcal R_{\rm rem}(u),
\]
where $\delta\hat C$ is the renormalized log-divergent constant and $\mathcal R_{\rm rem}$ is a $\mu_0$-a.s.\ finite (in fact polynomial-in-Wick-quantities) functional that involves third-order chaos terms and the ``$\Pi$'' tree expansions; its precise form is in \cite[\S 9]{Hairer14}.  Each Wick-polynomial $:\!u^k\!:$ is a well-defined element of the $k$-th Wiener chaos of $\mu_0$, an element of the negative-regularity Sobolev space $H^{-(k-1)/2-\kappa}$ ($\mu_0$-a.s.).

\textbf{Key identity.}  For $\psi\in C^\infty$,
\begin{align*}
:\!(u-\psi)^4\!: &= :\!u^4\!:\,-\,4\psi\,:\!u^3\!:\,+\,6\psi^2\,:\!u^2\!:\,-\,4\psi^3 u + \psi^4,\\
:\!(u-\psi)^2\!: &= :\!u^2\!:\,-\,2\psi u + \psi^2.
\end{align*}
These are honest equalities at the level of Wick-renormalized random variables on the same probability space (the renormalization constants $C_\varepsilon$ in $:\!(u-\psi)^k\!:$ and $:\!u^k\!:$ are the \emph{same}, because they depend only on the kernel $(m^2-\Delta)^{-1}$, not on the function being chaos-expanded).  Each term on the right is integrable against any test function, in particular against $1$ on $\mathbb T^3$:
\begin{itemize}\setlength\itemsep{2pt}
\item $\int\psi\,:\!u^3\!:\,$, $\int\psi^2\,:\!u^2\!:\,$, $\int\psi^3 u$ are all $\mu_0$-a.s.\ finite real-valued random variables (test the distribution-valued chaos $:\!u^k\!:$ against the smooth function $\psi^k$ or $\psi^{4-k}$).
\item $\int\psi^4$ is a deterministic finite constant.
\end{itemize}

Therefore
\begin{equation}\label{eq:diff_R}
\mathcal R(u)-\mathcal R(u-\psi) \;=\;\Phi(u,\psi)
\end{equation}
where
\[
\Phi(u,\psi)\;=\;\int_{\mathbb T^3}\!\bigl[\psi\,:\!u^3\!:\,-\,\tfrac32\psi^2\,:\!u^2\!:\,+\,\psi^3 u\,-\,\tfrac14\psi^4\bigr]dx \;+\; \tfrac32\delta\hat C\!\int(2\psi u-\psi^2)dx + \Phi_{\rm rem}(u,\psi),
\]
and $\Phi_{\rm rem}$ collects the analogous differences from $\mathcal R_{\rm rem}$, themselves well-defined $\mu_0$-a.s.\ functionals.  Each summand of $\Phi(u,\psi)$ is a $\mu_0$-a.s.\ finite real number, hence so is $\Phi(u,\psi)$.

\textbf{Step 4.}  Substituting \eqref{eq:diff_R} into \eqref{eq:RN_phi4}:
\[
\frac{dT_\psi^*\mu}{d\mu}(u)\;=\;e^{\Phi(u,\psi)}\cdot D_\psi(u).
\]
Both factors are strictly positive and $\mu_0$-a.s.\ finite, hence $\mu$-a.s.\ finite (since $\mu\ll\mu_0$).  Reversing the roles of $\mu$ and $T_\psi^*\mu$ (replace $\psi$ by $-\psi$), the same argument gives a strictly positive Radon--Nikodym derivative for $\mu$ w.r.t.\ $T_\psi^*\mu$.  Hence $\mu\sim T_\psi^*\mu$.

\medskip
\noindent\textbf{Integrability.}  $D_\psi\in L^1(\mu_0)$ by Cameron--Martin.  $e^{\Phi(u,\psi)}\in L^1(\mu)$ because: by Hairer's construction, $\mu$ has Gaussian-like tails for Wick polynomials of $u$ tested against smooth functions (each $\int\psi^k\,:\!u^j\!:\,dx$ is a $\mu$-integrable random variable with all polynomial moments; this is the $L^p$-integrability of Wick polynomials in the chaos).  The polynomial-in-$u$ exponent $\Phi(u,\psi)$ is dominated, modulo lower-order chaos, by the same quartic Wick polynomial that defines $\mathcal R$, giving the required exponential integrability via the standard $\Phi^4$ a-priori bounds.  This concludes the proof of equivalence and integrability.\qed

\section*{Remarks}

\begin{enumerate}\setlength\itemsep{4pt}
\item \textbf{Cameron--Martin space of $\mu$.}  The proof shows that the Cameron--Martin space of $\mu$ contains $C^\infty(\mathbb T^3)$.  In fact one can show it equals $H^1(\mathbb T^3)$ (the same as the Cameron--Martin space of $\mu_0$): any shift in $H^1$ gives an equivalent measure, while shifts outside $H^1$ are singular, since they are already singular at the level of the Gaussian $\mu_0$ (Cameron--Martin is sharp for Gaussian measures).
\item \textbf{Why smoothness of $\psi$ is more than needed.}  Only $\psi\in H^1$ is used.  The hypothesis $\psi\in C^\infty$ in the problem statement is therefore not optimal but more than sufficient.  The hypothesis ``$\psi\not\equiv 0$'' is irrelevant to the conclusion (the trivial shift is also equivalent — it is the identity).
\item \textbf{Extension.}  The same argument shows $T_\psi^*\mu\sim\mu$ for any $\psi$ in the Cameron--Martin space of $\mu_0$ on which the polynomial $:\!(u-\psi)^k\!:$ identities make sense, i.e., $\psi\in H^1\cap L^4(\mathbb T^3)=H^1$ (Sobolev embedding $H^1(\mathbb T^3)\hookrightarrow L^6$).
\end{enumerate}

\begin{thebibliography}{9}
\bibitem{Hairer14} M.~Hairer.  \emph{A theory of regularity structures}.  Invent.\ Math.\ \textbf{198} (2014), 269--504.
\bibitem{HairerMatetski18} M.~Hairer, K.~Matetski.  \emph{Discretisations of rough stochastic PDEs}.  Ann.\ Probab.\ \textbf{46} (2018).
\bibitem{BG20} N.~Barashkov, M.~Gubinelli.  \emph{A variational method for $\Phi^4_3$}.  Duke Math.\ J.\ \textbf{169} (2020), 3339--3415.
\bibitem{GH21} M.~Gubinelli, M.~Hofmanová.  \emph{A PDE construction of the Euclidean $\Phi^4_3$ quantum field theory}.  Comm.\ Math.\ Phys.\ \textbf{384} (2021), 1--75.
\bibitem{DPZ} G.~Da~Prato, J.~Zabczyk.  \emph{Stochastic Equations in Infinite Dimensions}, 2nd ed., CUP, 2014.
\end{thebibliography}

\end{document}
